\documentclass[../libro.tex]{subfiles}

\begin{document}

\ifSubfilesClassLoaded{\appendix\chapter{Afterstories}\clearpage}{}

\section{行列式的性质}
% \section{我要介绍行列式的哪些性质?}

% 这是一个好问题.

% 我初学行列式时, 教材讲了好几条性质.
% 我姑且用第一章的话译它们
% 这里,
设 \(a_1\), \(a_2\), \(\dots\), \(a_n\), \(x\), \(y\)
是 \(n+2\)~个 \(n \times 1\)~阵.
设 \(A = [a_1, a_2, \dots, a_n]\).
设 \(s\) 是一个数.
行列式有如下性质:

\vspace{2ex}

(1)
一个方阵与它的转置的行列式相等.
于是, 我不必同时讲行的性质与列的性质,
因为您总可用转置,
译列的性质为行的性质
(或译行的性质为列的性质).

(2)
若一个阵的某一列是二组数的和,
那么此阵的行列式等于二个阵的行列式的和,
而这二个阵, 除这一列以外, 全与原阵的对应的列一样.
用公式写, 就是
\begin{align*}
         & \det {[\dots, a_{j-1}, x + y, a_{j+1}, \dots]}
    \\
    = {} & \det {[\dots, a_{j-1}, x, a_{j+1}, \dots]}
    + \det {[\dots, a_{j-1}, y, a_{j+1}, \dots]},
\end{align*}
其中, 未写的列是不变的, 下同.

(3)
以一个数乘阵的一列后得到的阵的行列式%
等于以此数乘原阵的行列式.
用公式写, 就是
\begin{align*}
    \det {[\dots, a_{j-1}, sx, a_{j+1}, \dots]}
    = s \det {[\dots, a_{j-1}, x, a_{j+1}, \dots]}.
\end{align*}

(4)
若一个阵有二列相同, 则它的行列式为 \(0\).

(5)
若交换一个阵的二列, 则行列式变号.

(6)
若一个阵有二列成比例, 则它的行列式为 \(0\).
(利用性质~(3) (4) 即知.)

(7)
若加一个阵的一列的倍于另一列, 则行列式不变.
(利用性质~(2) (6) 即知.)

(8)
单位阵的行列式是 \(1\).

(9)
设 \(j\) 是不超过 \(n\) 的正整数.
则
\begin{align*}
    \det {(A)} = \sum_{i = 1}^{n}
    {(-1)^{i+j} [A]_{i,j} \det {(A(i|j))}}.
\end{align*}

\vspace{2ex}

不难看出, (2) (3) 的联合是多线性,
且 (4) 就是交错性.

% 我在第一章讲了行列式的 5~个性质.
% 用行列式与转置的关系可译行列式的关于阵的列的命题%
% 为关于行的命题.
% 规范性指出, 单位阵的行列式是 \(1\);
% 这是前面 7~条性质未指出的.
% 多线性就是 (2) (3) 的联合.
% 交错性就是 (4).
% 反称性就是 (5).
% 我没介绍 (6) (7); 它没介绍规范性:
% 我们二个都有好的未来.
% (这句话里包含三个分句,
% 前面二个是并列的,
% 末了儿一个是总起来说的,
% 就得在 ``规范性'' 之后用冒号,
% 表示并列的二个分句已完了,
% 还有总起来说的一个分句在后面.
% 这同样是提示的意思.
% 如此用冒号, 一些人还没有养成习惯.
% 可是, 这个用法可明白地表示分句与分句的关系,
% 应学会它.)

这些性质是有用的.
我们看几个例.

\begin{example}
    设 \(n\)~级阵 \(A\) 适合,
    当 \(1 \leq j < i \leq n\) 时,
    \([A]_{i,j} = 0\).
    通俗地,
    \begin{align*}
        A =
        \begin{bmatrix}
            [A]_{1,1} & [A]_{1,2} & \cdots & [A]_{1,n-1}   & [A]_{1,n}   \\
            0         & [A]_{2,2} & \cdots & [A]_{2,n-1}   & [A]_{2,n}   \\
            \vdots    & \vdots    & \ddots & \vdots        & \vdots      \\
            0         & 0         & \cdots & [A]_{n-1,n-1} & [A]_{n-1,n} \\
            0         & 0         & \cdots & 0             & [A]_{n,n}   \\
        \end{bmatrix}.
    \end{align*}
    我们计算 \(\det {(A)}\).

    按列~\(1\) 展开, 有
    \begin{align*}
        \det {(A)}
        = {} &
        \sum_{i=1}^{n}
        {(-1)^{i+1} [A]_{i,1} \det {(A(i|1))}}
        \\
        = {} &
        (-1)^{1+1} [A]_{1,1} \det {(A(1|1))}
        + \sum_{i=2}^{n}
        {(-1)^{i+1} [A]_{i,1} \det {(A(i|1))}}
        \\
        = {} &
        [A]_{1,1} \det {(A(1|1))}
        + \sum_{i=2}^{n}
        {(-1)^{i+1}\, 0 \det {(A(i|1))}}
        \\
        = {} &
        [A]_{1,1} \det {(A(1|1))}.
    \end{align*}

    不难看出, 当 \(1 \leq j < i \leq n-1\) 时,
    也有 \([A(1|1)]_{i,j} = 0\).
    于是, 类似地,
    \begin{align*}
        \det {(A(1|1))}
        = [A(1|1)]_{1,1} \det {((A(1|1))(1|1))}
            = [A]_{2,2} \det {(A({1,2}|{1,2}))}.
    \end{align*}
    故
    \begin{align*}
        \det {(A)}
        = [A]_{1,1} [A]_{2,2} \det {(A({1,2}|{1,2}))}.
    \end{align*}

    ……

    最后, 我们得
    \begin{align*}
        \det {(A)}
        = {} &
        [A]_{1,1} [A]_{2,2} \dots [A]_{n-1,n-1}
        \det {(A({1,2,\dots,n-1}|{1,2,\dots,n-1}))}
        \\
        = {} &
        [A]_{1,1} [A]_{2,2} \dots [A]_{n-1,n-1} [A]_{n,n}
        \\
        = {} &
        [A]_{1,1} [A]_{2,2} \dots [A]_{n,n}.
    \end{align*}
\end{example}

\begin{example}
    设 \(n\)~级阵 \(A\) 适合,
    当 \(1 \leq i < j \leq n\) 时,
    \([A]_{i,j} = 0\).
    通俗地,
    \begin{align*}
        A =
        \begin{bmatrix}
            [A]_{1,1}   & 0           & \cdots & 0             & 0         \\
            [A]_{2,1}   & [A]_{2,2}   & \cdots & 0             & 0         \\
            \vdots      & \vdots      & \ddots & \vdots        & \vdots    \\
            [A]_{n-1,1} & [A]_{n-1,2} & \cdots & [A]_{n-1,n-1} & 0         \\
            [A]_{n,1}   & [A]_{n,2}   & \cdots & [A]_{n,n-1}   & [A]_{n,n} \\
        \end{bmatrix}.
    \end{align*}
    我们计算 \(\det {(A)}\).

    不难看出, \(A\) 的转置 \(A^{\mathrm{T}}\) 适合,
    当 \(1 \leq j < i \leq n\) 时,
    \([A^{\mathrm{T}}]_{i,j} = [A]_{j,i} = 0\).
    由性质~(1) 与上例的结果,
    \begin{align*}
        \det {(A)}
        = {} &
        \det {(A^{\mathrm{T}})}
        \\
        = {} &
        [A^{\mathrm{T}}]_{1,1} [A^{\mathrm{T}}]_{2,2}
        \dots [A^{\mathrm{T}}]_{n,n}
        \\
        = {} &
        [A]_{1,1} [A]_{2,2} \dots [A]_{n,n}.
    \end{align*}
\end{example}

\begin{example}
    运用性质~(7), 可作出一些 \(0\), 使计算简单.
    比如,
    设
    \begin{align*}
        A = \begin{bmatrix}
                4 & 9 & 2 \\
                3 & 5 & 7 \\
                8 & 1 & 6 \\
            \end{bmatrix}.
    \end{align*}
    则
    \begin{align*}
        \det {(A)}
        = {} &
        \det {
            \begin{bmatrix}
                4 & 9 & 2 + 9 \\
                3 & 5 & 7 + 5 \\
                8 & 1 & 6 + 1 \\
            \end{bmatrix}
        }
        \\
        = {} &
        \det {
            \begin{bmatrix}
                4 & 9 & 2 + 9 + 4 \\
                3 & 5 & 7 + 5 + 3 \\
                8 & 1 & 6 + 1 + 8 \\
            \end{bmatrix}
        }
        \\
        = {} &
        \det {
            \begin{bmatrix}
                4 & 9 & 15 \\
                3 & 5 & 15 \\
                8 & 1 & 15 \\
            \end{bmatrix}
        }
        \\
        = {} &
        \det {
            \begin{bmatrix}
                4 & 9 - 15 \cdot \frac{1}{15} & 15 \\
                3 & 5 - 15 \cdot \frac{1}{15} & 15 \\
                8 & 1 - 15 \cdot \frac{1}{15} & 15 \\
            \end{bmatrix}
        }
        \\
        = {} &
        \det {
            \begin{bmatrix}
                4 - 15 \cdot \frac{8}{15} & 8 & 15 \\
                3 - 15 \cdot \frac{8}{15} & 4 & 15 \\
                8 - 15 \cdot \frac{8}{15} & 0 & 15 \\
            \end{bmatrix}
        }
        \\
        = {} &
        \det {
            \begin{bmatrix}
                -4 & 8 & 15 \\
                -5 & 4 & 15 \\
                0  & 0 & 15 \\
            \end{bmatrix}
        }
        \\
        = {} &
        \det {
            \begin{bmatrix}
                -4 + 8 \cdot \frac{5}{4} & 8 & 15 \\
                -5 + 4 \cdot \frac{5}{4} & 4 & 15 \\
                0 + 0 \cdot \frac{5}{4}  & 0 & 15 \\
            \end{bmatrix}
        }
        \\
        = {} &
        \det {
            \begin{bmatrix}
                6 & 8 & 15 \\
                0 & 4 & 15 \\
                0 & 0 & 15 \\
            \end{bmatrix}
        }
        \\
        = {} &
        6 \cdot 4 \cdot 15
        \\
        = {} &
        360.
    \end{align*}
    我们用 \(3\)~级阵的行列式的公式验证结果:
    \begin{align*}
        \det {(A)}
        = {} &
        4 \cdot 5 \cdot 6
        + 3 \cdot 1 \cdot 2
        + 8 \cdot 9 \cdot 7
        - 4 \cdot 1 \cdot 7
        - 3 \cdot 9 \cdot 6
        - 8 \cdot 5 \cdot 2
        \\
        = {} &
        120 + 6 + 504 - 28 - 162 - 80
        \\
        = {} & 360.
    \end{align*}
\end{example}

\begin{example}
    设
    \begin{align*}
        A = \begin{bmatrix}
                16 & 3  & 2  & 13 \\
                5  & 10 & 11 & 8  \\
                9  & 6  & 7  & 12 \\
                4  & 15 & 14 & 1  \\
            \end{bmatrix}.
    \end{align*}
    则
    \begin{align*}
        \det {(A)}
        = {} &
        \det {
            \begin{bmatrix}
                16 & 3 - 2   & 2  & 13 \\
                5  & 10 - 11 & 11 & 8  \\
                9  & 6 - 7   & 7  & 12 \\
                4  & 15 - 14 & 14 & 1  \\
            \end{bmatrix}
        }
        \\
        = {} &
        \det {
            \begin{bmatrix}
                16 - 13 & 3 - 2   & 2  & 13 \\
                5 - 8   & 10 - 11 & 11 & 8  \\
                9 - 12  & 6 - 7   & 7  & 12 \\
                4 - 1   & 15 - 14 & 14 & 1  \\
            \end{bmatrix}
        }
        \\
        = {} &
        \det {
            \begin{bmatrix}
                3  & 1  & 2  & 13 \\
                -3 & -1 & 11 & 8  \\
                -3 & -1 & 7  & 12 \\
                3  & 1  & 14 & 1  \\
            \end{bmatrix}
        }.
    \end{align*}
    利用性质~(6) 可知, 最后一个阵的行列式为 \(0\).
    所以, \(A\)~的行列式也是 \(0\).

    我们当然也可用
    \(4\)~级阵的行列式的公式验证结果.
    不过, 这是复杂的.
    首先, 根据定义,
    \begin{align*}
        %  &
        \det {(A)}
        %  \\
        = {} &
        \hphantom{{} + {}}
        (-1)^{1+1} [A]_{1,1} \det {(A(1|1))}
        + (-1)^{2+1} [A]_{2,1} \det {(A(2|1))}
        \\
             &
        + (-1)^{3+1} [A]_{3,1} \det {(A(3|1))}
        + (-1)^{4+1} [A]_{4,1} \det {(A(4|1))}.
    \end{align*}
    可算出
    \begin{align*}
         & \det {(A(1|1))}
        = \det {\begin{bmatrix}
                        10 & 11 & 8  \\
                        6  & 7  & 12 \\
                        15 & 14 & 1  \\
                    \end{bmatrix}}
        = 136;             \\
         & \det {(A(2|1))}
        = \det {\begin{bmatrix}
                        3  & 2  & 13 \\
                        6  & 7  & 12 \\
                        15 & 14 & 1  \\
                    \end{bmatrix}}
        = -408;            \\
         & \det {(A(3|1))}
        = \det {\begin{bmatrix}
                        3  & 2  & 13 \\
                        10 & 11 & 8  \\
                        15 & 14 & 1  \\
                    \end{bmatrix}}
        = -408;            \\
         & \det {(A(4|1))}
        = \det {\begin{bmatrix}
                        3  & 2  & 13 \\
                        10 & 11 & 8  \\
                        6  & 7  & 12 \\
                    \end{bmatrix}}
        = 136.
    \end{align*}
    所以
    \begin{align*}
        \det {(A)}
        = 16 \cdot 136 - 5 \cdot (-408)
        + 9 \cdot (-408) - 4 \cdot 136
        = 0.
    \end{align*}
\end{example}

\begin{example}
    设 \(x\), \(y\) 是数.
    作 \(n\)~级阵 \(A\) 如下:
    \begin{align*}
        [A]_{i,j} =
        \begin{cases}
            x, & i = j;    \\
            y, & i \neq j.
        \end{cases}
    \end{align*}
    通俗地,
    \begin{align*}
        A =
        \begin{bmatrix}
            x      & y      & \cdots & y      & y      \\
            y      & x      & \cdots & y      & y      \\
            \vdots & \vdots & {}     & \vdots & \vdots \\
            y      & y      & \cdots & x      & y      \\
            y      & y      & \cdots & y      & x      \\
        \end{bmatrix}.
    \end{align*}
    我们计算 \(\det {(A)}\).

    加 \(A\)~的%
    列~\(1\), \(2\), \(\dots\), \(n-1\) 于列~\(n\),
    得 \(n\)~级阵
    \begin{align*}
        A_1 =
        \begin{bmatrix}
            x      & y      & \cdots & y      & x + (n-1)y \\
            y      & x      & \cdots & y      & x + (n-1)y \\
            \vdots & \vdots & {}     & \vdots & \vdots     \\
            y      & y      & \cdots & x      & x + (n-1)y \\
            y      & y      & \cdots & y      & x + (n-1)y \\
        \end{bmatrix}.
    \end{align*}
    注意, \(A_1\) 的前 \(n-1\)~列%
    与 \(A\) 的前 \(n-1\)~列一样,
    但 \(A_1\) 的列~\(n\) 的元全为 \(x + (n-1)y\).
    用 \(n-1\)~次性质~(7), 有
    \(\det {(A)} = \det {(A_1)}\).

    作 \(n\)~级阵
    \begin{align*}
        A_2 =
        \begin{bmatrix}
            x      & y      & \cdots & y      & 1      \\
            y      & x      & \cdots & y      & 1      \\
            \vdots & \vdots & {}     & \vdots & \vdots \\
            y      & y      & \cdots & x      & 1      \\
            y      & y      & \cdots & y      & 1      \\
        \end{bmatrix}.
    \end{align*}
    注意, \(A_2\) 的前 \(n-1\)~列%
    与 \(A_1\) 的前 \(n-1\)~列一样,
    但 \(A_2\) 的列~\(n\) 的元全为 \(1\).
    用性质~(3), 有
    \(\det {(A)} = \det {(A_1)}
    = (x + (n-1)y) \det {(A_2)}\).

    加 \(A_2\)~的%
    列~\(n\) 的 \(-y\)~倍于列~\(1\),
    列~\(n\) 的 \(-y\)~倍于列~\(2\),
    ……
    列~\(n\) 的 \(-y\)~倍于列~\(n-1\),
    得 \(n\)~级阵
    \begin{align*}
        A_3 =
        \begin{bmatrix}
            x-y    & 0      & \cdots & 0      & 1      \\
            0      & x-y    & \cdots & 0      & 1      \\
            \vdots & \vdots & {}     & \vdots & \vdots \\
            0      & 0      & \cdots & x-y    & 1      \\
            0      & 0      & \cdots & 0      & 1      \\
        \end{bmatrix}.
    \end{align*}
    用 \(n-1\)~次性质~(7), 有
    \(\det {(A)}
    = (x + (n-1)y) \det {(A_2)}
    = (x + (n-1)y) \det {(A_3)}\).
    注意,
    \(1 \leq j < i \leq n\) 时,
    有 \([A_3]_{i,j} = 0\).
    故
    \begin{align*}
        \det {(A_3)} = (x-y)^{n-1} \cdot 1 = (x-y)^{n-1}.
    \end{align*}
    故
    \begin{align*}
        \det {(A)}
        = (x + (n-1)y) \det {(A_3)}
        = (x + (n-1)y) (x-y)^{n-1}.
    \end{align*}
\end{example}

\begin{example}[Vandermonde \pinjino{fandemonde} 阵的行列式]
    设 \(x_1\), \(x_2\), \(\dots\), \(x_n\) 是数.
    作 \(n\)~级 Vandermonde 阵
    \(V(x_1, x_2, \dots, x_n)\) 如下:
    \begin{align*}
        [V(x_1, x_2, \dots, x_n)]_{i,j}
        = x_i^{j-1}.
    \end{align*}
    通俗地,
    \begin{align*}
        V(x_1, x_2, \dots, x_n)
        =
        \begin{bmatrix}
            1      & x_1     & x_1^2     & \cdots & x_1^{n-2}     & x_1^{n-1}     \\
            1      & x_2     & x_2^2     & \cdots & x_2^{n-2}     & x_2^{n-1}     \\
            1      & x_3     & x_3^2     & \cdots & x_3^{n-2}     & x_3^{n-1}     \\
            \vdots & \vdots  & \vdots    & {}     & \vdots        & \vdots        \\
            1      & x_{n-1} & x_{n-1}^2 & \cdots & x_{n-1}^{n-2} & x_{n-1}^{n-1} \\
            1      & x_n     & x_n^2     & \cdots & x_n^{n-2}     & x_n^{n-1}     \\
        \end{bmatrix}.
    \end{align*}
    我们计算 \(\det {(V(x_1, x_2, \dots, x_n))}\).

    加 \(V(x_1, x_2, \dots, x_n)\)~的%
    列~\(n-1\) 的 \(-x_n\) 倍于列~\(n\),
    列~\(n-2\) 的 \(-x_n\) 倍于列~\(n-1\),
    ……
    列~\(1\) 的 \(-x_n\) 倍于列~\(2\),
    得 \(n\)~级阵
    \begin{align*}
             & A
        \\
        = {} &
        \begin{bmatrix}
            1      & x_1 - x_n     & x_1^2 - x_1 x_n         & \cdots & x_1^{n-2} - x_1^{n-3} x_n         & x_1^{n-1} - x_1^{n-2} x_n         \\
            1      & x_2 - x_n     & x_2^2 - x_2 x_n         & \cdots & x_2^{n-2} - x_2^{n-3} x_n         & x_2^{n-1} - x_2^{n-2} x_n         \\
            1      & x_3 - x_n     & x_3^2 - x_3 x_n         & \cdots & x_3^{n-2} - x_3^{n-3} x_n         & x_3^{n-1} - x_3^{n-2} x_n         \\
            \vdots & \vdots        & \vdots                  & {}     & \vdots                            & \vdots                            \\
            1      & x_{n-1} - x_n & x_{n-1}^2 - x_{n-1} x_n & \cdots & x_{n-1}^{n-2} - x_{n-1}^{n-3} x_n & x_{n-1}^{n-1} - x_{n-1}^{n-2} x_n \\
            1      & 0             & 0                       & \cdots & 0                                 & 0                                 \\
        \end{bmatrix}
        \\
        = {} &
        \begin{bmatrix}
            1      & x_1 - x_n     & x_1 (x_1 - x_n)         & \cdots & x_1^{n-3} (x_1 - x_n)         & x_1^{n-2} (x_1 - x_n)         \\
            1      & x_2 - x_n     & x_2 (x_2 - x_n)         & \cdots & x_2^{n-3} (x_2 - x_n)         & x_2^{n-2} (x_2 - x_n)         \\
            1      & x_3 - x_n     & x_3 (x_3 - x_n)         & \cdots & x_3^{n-3} (x_3 - x_n)         & x_3^{n-2} (x_3 - x_n)         \\
            \vdots & \vdots        & \vdots                  & {}     & \vdots                        & \vdots                        \\
            1      & x_{n-1} - x_n & x_{n-1} (x_{n-1} - x_n) & \cdots & x_{n-1}^{n-3} (x_{n-1} - x_n) & x_{n-1}^{n-2} (x_{n-1} - x_n) \\
            1      & 0             & 0                       & \cdots & 0                             & 0                             \\
        \end{bmatrix};
    \end{align*}
    具体地,
    \begin{align*}
        [A]_{i,j}
        =
        \begin{cases}
            1,
             & j = 1; \\
            x_i^{j-1} - x_i^{j-2} x_n
            = x_i^{j-2} (x_i - x_n),
             & j > 2.
        \end{cases}
    \end{align*}
    由性质~(7), 有
    \begin{align*}
        \det {(V(x_1, x_2, \dots, x_n))} = \det {(A)}.
    \end{align*}

    按行~\(n\) 展开, 有
    \begin{align*}
        \det {(V(x_1, x_2, \dots, x_n))}
        = \det {(A)}
        = (-1)^{n+1} \det {(B)},
    \end{align*}
    其中, \(B = A(n|1)\),
    且 \([B]_{i,j} = x_i^{j-1} (x_i - x_n)\);
    通俗地,
    \begin{align*}
        B =
        \begin{bmatrix}
            x_1 - x_n     & x_1 (x_1 - x_n)         & \cdots & x_1^{n-3} (x_1 - x_n)         & x_1^{n-2} (x_1 - x_n)         \\
            x_2 - x_n     & x_2 (x_2 - x_n)         & \cdots & x_2^{n-3} (x_2 - x_n)         & x_2^{n-2} (x_2 - x_n)         \\
            x_3 - x_n     & x_3 (x_3 - x_n)         & \cdots & x_3^{n-3} (x_3 - x_n)         & x_3^{n-2} (x_3 - x_n)         \\
            \vdots        & \vdots                  & {}     & \vdots                        & \vdots                        \\
            x_{n-1} - x_n & x_{n-1} (x_{n-1} - x_n) & \cdots & x_{n-1}^{n-3} (x_{n-1} - x_n) & x_{n-1}^{n-2} (x_{n-1} - x_n) \\
        \end{bmatrix}.
    \end{align*}

    用 \(n-1\)~次性质~(3) (关于行), 有
    \begin{align*}
        \det {(B)}
        = {} &
        (x_1 - x_n) (x_2- x_n) \dots (x_{n-1} - x_n)
        \det {(V(x_1, x_2, \dots, x_{n-1}))}
        \\
        = {} &
        (-1)^{n-1} (x_n - x_1) (x_n - x_2) \dots (x_n - x_{n-1})
        \det {(V(x_1, x_2, \dots, x_{n-1}))}.
    \end{align*}
    故
    \begin{align*}
             &
        \det {(V(x_1, x_2, \dots, x_n))}
        \\
        = {} &
        (-1)^{n+1} \det {(B)}
        \\
        = {} &
        (-1)^{n+1}
        (-1)^{n-1} (x_n - x_1) (x_n - x_2) \dots (x_n - x_{n-1})
        \det {(V(x_1, x_2, \dots, x_{n-1}))}
        \\
        = {} &
        (x_n - x_1) (x_n - x_2) \dots (x_n - x_{n-1})
        \det {(V(x_1, x_2, \dots, x_{n-1}))}.
    \end{align*}

    类似地,
    \begin{align*}
             &
        \det {(V(x_1, x_2, \dots, x_{n-1}))}
        \\
        = {} &
        (x_{n-1} - x_1) (x_{n-1} - x_2) \dots (x_{n-1} - x_{n-2})
        \det {(V(x_1, x_2, \dots, x_{n-2}))}.
    \end{align*}
    故
    \begin{align*}
        \det {(V(x_1, x_2, \dots, x_n))}
        = {} & \hphantom{\cdot\,\,}
        (x_n - x_1) (x_n - x_2) \dots (x_n - x_{n-1})
        \\
             & \cdot
        (x_{n-1} - x_1) (x_{n-1} - x_2) \dots (x_{n-1} - x_{n-2})
        \\
             & \cdot
        \det {(V(x_1, x_2, \dots, x_{n-2}))}.
    \end{align*}

    ……

    最后, 我们有
    \begin{align*}
        \det {(V(x_1, x_2, \dots, x_n))}
        = {} & \hphantom{\cdot\,\,}
        (x_n - x_1) (x_n - x_2) \dots (x_n - x_{n-1})
        \\
             & \cdot
        (x_{n-1} - x_1) (x_{n-1} - x_2) \dots (x_{n-1} - x_{n-2})
        \\
             & \cdot
        \dots
        \\
             & \cdot
        (x_3 - x_1) (x_3 - x_2)
        \\
             & \cdot
        (x_2 - x_1)
        \\
             & \cdot
        \det {(V(x_1))}
        \\
        = {} & \hphantom{\cdot\,\,}
        (x_n - x_1) (x_n - x_2) \dots (x_n - x_{n-1})
        \\
             & \cdot
        (x_{n-1} - x_1) (x_{n-1} - x_2) \dots (x_{n-1} - x_{n-2})
        \\
             & \cdot
        \dots
        \\
             & \cdot
        (x_3 - x_1) (x_3 - x_2)
        \\
             & \cdot
        (x_2 - x_1)
        \\
             & \cdot
        1
        \\
        = {} &
        \prod_{2 \leq v \leq n}
        {
            \prod_{1 \leq u \leq v-1} {(x_v - x_u)}
        }
        \\
        = {} &
        \prod_{1 \leq u < v \leq n} {(x_v - x_u)}.
    \end{align*}

    不难看出,
    若 \(V(x_1, x_2, \dots, x_n)\) 的行列式非零,
    则 \(x_1\), \(x_2\), \(\dots\), \(x_n\) 互不相同;
    反过来,
    若 \(x_1\), \(x_2\), \(\dots\), \(x_n\) 互不相同,
    则 \(V(x_1, x_2, \dots, x_n)\) 的行列式非零.
\end{example}

\begin{example}
    设 \(a_0\), \(a_1\), \(a_2\), \(\dots\), \(a_n\) 是 \(n+1\)~个数.
    作 \(n\)~级阵 \(P(x; a_0; a_1, a_2, \dots, a_n)\) 如下:
    \begin{align*}
        [P(x; a_0; a_1, a_2, \dots, a_n)]_{i,j}
        = \begin{cases}
              x,           & 1 \leq i = j \leq n-1;   \\
              -1,          & 1 \leq i-1 = j \leq n-1; \\
              a_{n-i+1},   & 1 \leq i \leq n-1 = j-1; \\
              a_0 x + a_1, & i = n = j;               \\
              0,           & \text{其他}.
          \end{cases}
    \end{align*}
    通俗地,
    \begin{align*}
        P(x; a_0; a_1, a_2, \dots, a_n)
        = \begin{bmatrix}
              x      & 0      & \cdots & 0      & 0      & a_n         \\
              -1     & x      & \cdots & 0      & 0      & a_{n-1}     \\
              0      & -1     & \cdots & 0      & 0      & a_{n-2}     \\
              \vdots & \vdots & {}     & \vdots & \vdots & \vdots      \\
              0      & 0      & \cdots & -1     & x      & a_2         \\
              0      & 0      & \cdots & 0      & -1     & a_0 x + a_1 \\
          \end{bmatrix}.
    \end{align*}
    我们计算 \(\det {(P(x; a_0; a_1, a_2, \dots, a_n))}\).

    注意, 当 \(1 < j < n\) 时,
    \([P(x; a_0; a_1, a_2, \dots, a_n)]_{1,j} = 0\).
    于是, 按行~\(1\) 展开, 有
    \begin{align*}
             &
        \det {(P(x; a_0; a_1, a_2, \dots, a_n))}
        \\
        = {} &
        \hphantom{{} + {}}
        (-1)^{1+1} [P(x; a_0; a_1, a_2, \dots, a_n)]_{1,1}
        \det {((P(x; a_0; a_1, a_2, \dots, a_n))(1|1))}
        \\
             &
        + (-1)^{1+n} [P(x; a_0; a_1, a_2, \dots, a_n)]_{1,n}
        \det {((P(x; a_0; a_1, a_2, \dots, a_n))(1|n))}.
    \end{align*}
    不难写出,
    \begin{align*}
        [P(x; a_0; a_1, a_2, \dots, a_n)]_{1,1} = x,
        \quad
        [P(x; a_0; a_1, a_2, \dots, a_n)]_{1,n} = a_n.
    \end{align*}
    不难写出,
    \begin{align*}
        (P(x; a_0; a_1, a_2, \dots, a_n))(1|1)
        = {} &
        \begin{bmatrix}
            x      & \cdots & 0      & 0      & a_{n-1}     \\
            -1     & \cdots & 0      & 0      & a_{n-2}     \\
            \vdots & {}     & \vdots & \vdots & \vdots      \\
            0      & \cdots & -1     & x      & a_2         \\
            0      & \cdots & 0      & -1     & a_0 x + a_1 \\
        \end{bmatrix}
        \\
        = {} &
        P(x; a_0; a_1, a_2, \dots, a_{n-1}).
    \end{align*}
    故
    \begin{align*}
        \det {((P(x; a_0; a_1, a_2, \dots, a_n))(1|1))}
        = \det {(P(x; a_0; a_1, a_2, \dots, a_{n-1}))}.
    \end{align*}
    不难写出,
    \begin{align*}
        (P(x; a_0; a_1, a_2, \dots, a_n))(1|n)
        = \begin{bmatrix}
              -1     & x      & \cdots & 0      & 0      \\
              0      & -1     & \cdots & 0      & 0      \\
              \vdots & \vdots & {}     & \vdots & \vdots \\
              0      & 0      & \cdots & -1     & x      \\
              0      & 0      & \cdots & 0      & -1     \\
          \end{bmatrix};
    \end{align*}
    具体地,
    \begin{align*}
        [(P(x; a_0; a_1, a_2, \dots, a_n))(1|n)]_{i,j}
        = \begin{cases}
              -1, & 1 \leq i = j \leq n-1;     \\
              x,  & 1 \leq i = j - 1 \leq n-2; \\
              0,  & \text{其他}.
          \end{cases}
    \end{align*}
    故 \(1 \leq j < i \leq n-1\) 时, 有
    \([(P(x; a_0; a_1, a_2, \dots, a_n))(1|n)]_{i,j} = 0\).
    则
    \begin{align*}
        \det {((P(x; a_0; a_1, a_2, \dots, a_n))(1|n))}
        = (-1)^{n-1}.
    \end{align*}

    综上,
    \begin{align*}
             &
        \det {(P(x; a_0; a_1, a_2, \dots, a_n))}
        \\
        = {} &
        \hphantom{{} + {}}
        (-1)^{1+1} [P(x; a_0; a_1, a_2, \dots, a_n)]_{1,1}
        \det {((P(x; a_0; a_1, a_2, \dots, a_n))(1|1))}
        \\
             &
        + (-1)^{1+n} [P(x; a_0; a_1, a_2, \dots, a_n)]_{1,n}
        \det {((P(x; a_0; a_1, a_2, \dots, a_n))(1|n))}
        \\
        = {} &
        x \det {(P(x; a_0; a_1, a_2, \dots, a_{n-1}))}
        + (-1)^{1+n} a_n (-1)^{n-1}
        \\
        = {} &
        x \det {(P(x; a_0; a_1, a_2, \dots, a_{n-1}))} + a_n.
    \end{align*}

    类似地,
    \begin{align*}
        \det {(P(x; a_0; a_1, a_2, \dots, a_{n-1}))}
        = x \det {(P(x; a_0; a_1, a_2, \dots, a_{n-2}))} + a_{n-1}.
    \end{align*}
    故
    \begin{align*}
        \det {(P(x; a_0; a_1, a_2, \dots, a_n))}
        = {} &
        x \det {(P(x; a_0; a_1, a_2, \dots, a_{n-1}))} + a_n
        \\
        = {} &
        x (x \det {(P(x; a_0; a_1, a_2, \dots, a_{n-2}))} + a_{n-1}) + a_n
        \\
        = {} &
        x^2 \det {(P(x; a_0; a_1, a_2, \dots, a_{n-2}))}
        + a_{n-1} x + a_n.
    \end{align*}

    ……

    最后, 我们有
    \begin{align*}
             &
        \det {(P(x; a_0; a_1, a_2, \dots, a_n))}
        \\
        = {} &
        x \det {(P(x; a_0; a_1, a_2, \dots, a_{n-1}))} + a_n
        \\
        = {} &
        x^2 \det {(P(x; a_0; a_1, a_2, \dots, a_{n-2}))} + a_{n-1} x + a_n
        \\
        = {} &
        \dots
        \\
        = {} &
        x^{n-1} \det {(P(x; a_0; a_1))} + a_2 x^{n-2} + \dots + a_{n-1} x + a_n
        \\
        = {} &
        x^{n-1} (a_0 x + a_1) + a_2 x^{n-2} + \dots + a_{n-1} x + a_n
        \\
        = {} &
        a_0 x^n + a_1 x^{n-1} + a_2 x^{n-2} + \dots + a_{n-1} x + a_n.
    \end{align*}
\end{example}

虽然我写了多的例,
我并不想在本书具体地介绍%
如何计算一个阵的行列式.
% 那或许不是我的事.
% 我的目标只是带您入门行列式;
这些例的目的是使您更好地理解行列式的性质.
% 若您对计算一个阵的行列式的方法感兴趣,
若您想知道计算一个阵的行列式的更多的方法,
您可以见一些线性代数教材,
或者找相关的文献.

\section{排列}

% 一般地, 在以组合定义定义行列式的教材里,
% 排列 (或置换) 的一些基本的性质是必要的,
% 因为它们对论证行列式的性质有用
% (通俗地, 这些教材用排列研究行列式).
% 我想在本节用行列式研究排列
% (通俗地, 我要反过来作).

本节, 我们讨论排列.

研究排列时, 有一个行为是重要的.

\begin{definition}[对换]
    设 \(a_1\), \(a_2\), \(\dots\), \(a_n\) 是一个排列.
    设 \(i\), \(j\) 是不超过 \(n\) 的二个不等的正整数.
    交换此排列的第~\(i\) 个文字与第~\(j\) 个文字
    (也就是, 交换文字 \(a_i\), \(a_j\)),
    且不改变其他的文字,
    则, 我们得到一个新的排列
    \(b_1\), \(b_2\), \(\dots\), \(b_n\),
    其中,
    \begin{align*}
        b_k
        = \begin{cases}
              a_j, & k = i;     \\
              a_i, & k = j;     \\
              a_k, & \text{其他}.
          \end{cases}
    \end{align*}
    我们说, 变 \(a_1\), \(a_2\), \(\dots\), \(a_n\)
    为 \(b_1\), \(b_2\), \(\dots\), \(b_n\)
    的行为是一次\emph{对换}.
    特别地, 若 \(j = i + 1\) 或
    \(j = i - 1\)
    (通俗地,
    \(a_i\), \(a_j\), 或 \(a_j\), \(a_i\), 在
    \(a_1\), \(a_2\), \(\dots\), \(a_n\) 里,
    是相邻的二个文字),
    我们说, 这一次对换是一次\emph{相邻对换}.
\end{definition}

\begin{example}
    考虑排列~I: \(3\), \(2\), \(5\), \(1\), \(4\).
    交换 \(3\) 与 \(1\) 的位置,
    得排列~II: \(1\), \(2\), \(5\), \(3\), \(4\).
    我们说, 变 I 为 II 的行为是一次对换.

    还是考虑排列~I.
    交换 \(1\) 与 \(4\) 的位置,
    得排列~III: \(3\), \(2\), \(5\), \(4\), \(1\).
    变 I 为 III 的行为当然也是一次对换.
    不过, 因为 \(1\), \(4\) 在 I 里是%
    相邻的二个文字,
    故我们也可说,
    变 I 为 III 的行为是一次相邻对换.

    注意,
    一次对换的结果可以跟多次相邻对换的结果一样.
    比如, 为了变 I 为 II,
    我们可以这么作:

    交换 \(3\) 与 \(2\) 的位置,
    得排列~IV: \(2\), \(3\), \(5\), \(1\), \(4\);

    交换 \(3\) 与 \(5\) 的位置,
    得排列~V: \(2\), \(5\), \(3\), \(1\), \(4\);

    交换 \(3\) 与 \(1\) 的位置,
    得排列~VI: \(2\), \(5\), \(1\), \(3\), \(4\);

    交换 \(5\) 与 \(1\) 的位置,
    得排列~VII: \(2\), \(1\), \(5\), \(3\), \(4\);

    交换 \(2\) 与 \(1\) 的位置,
    得排列~VIII: \(1\), \(2\), \(5\), \(3\), \(4\).

    可以看到,
    我们利用 \(5\)~次相邻对换实现了%
    一次 (不是相邻的) 对换.
    注意, 这是一个奇数.

    其实, 不难看出这么一件事.
    设 \(a_1\), \(a_2\), \(\dots\), \(a_n\) 是一个排列,
    \(i\), \(j\) 是不超过 \(n\) 的二个不等的正整数,
    且 \(i < j\).
    则, 我们可作 \((j - i) + (j - 1 - i)
    = 2(j - i) - 1\)~次相邻对换%
    以交换此排列的第~\(i\) 个文字与第~\(j\) 个文字.
\end{example}

相邻对换与逆序数有如下关系.

\begin{theorem}
    设 \(a_1\), \(a_2\), \(\dots\), \(a_n\)
    是互不相同的 \(n\)~个整数.
    设排列~I: \(a_1\), \(a_2\), \(\dots\), \(a_n\)
    的逆序数 \(\tau (a_1, a_2, \dots, a_n) = T\).
    那么, 我们可作 \(T\)~次相邻对换%
    变排列~I 为它的自然排列.
\end{theorem}

% 我宣布一件事.
% 从现在开始, 我所说的 ``文字列'' 的文字全为整数;
% 自然地, 我所说的 ``排列'' 的文字也全为整数.

\begin{proof}
    我们先设 \(a_1\), \(a_2\), \(\dots\), \(a_n\)
    是 \(1\), \(2\), \(\dots\), \(n\) 的排列.
    作命题 \(P(n)\):
    \begin{quotation}
        对 \(1\), \(2\), \(\dots\), \(n\)
        的任何排列 \(a_1\), \(a_2\), \(\dots\), \(a_n\),
        我们可作 \(\tau (a_1, a_2, \dots, a_n)\)~次相邻对换%
        变它为自然排列 \(1\), \(2\), \(\dots\), \(n\).
    \end{quotation}
    我们用数学归纳法证明,
    对任何正整数 \(n\), \(P(n)\) 成立.

    \(P(1)\) 是对的.
    \(1\) 只有 \(1\)~个排列: \(1\).
    我们不必作任何对换, 它已是自然排列.
    排列 \(1\)~的逆序数自然是 \(0\).

    现在, 我们设 \(P(m-1)\) 是对的.
    我们证 \(P(m)\) 也是对的.
    任取 \(1\), \(2\), \(\dots\), \(m\) 的一个排列
    \(a_1\), \(a_2\), \(\dots\), \(a_m\).
    那么, 恰存在一个不超过 \(m\) 的正整数 \(k\),
    使 \(a_k = m\).
    交换 \(m\) 与 \(a_{k+1}\),
    再交换 \(m\) 与 \(a_{k+2}\),
    ……
    再交换 \(m\) 与 \(a_m\),
    从而可变 \(a_1\), \(a_2\), \(\dots\), \(a_m\)
    为 \(a_1\), \(\dots\), \(a_{k-1}\),
    \(a_{k+1}\), \(\dots\), \(a_m\), \(m\).
    这 \(m - k\)~次相邻对换%
    使 \(m\) 是最后一个数
    % 使 \(m\) 在最末的位置
    (若 \(a_m = m\), 则不必作对换了,
    故此关系仍成立),
    且新排列的前 \(m-1\)~个整数
    \(a_1\), \(\dots\), \(a_{k-1}\),
    \(a_{k+1}\), \(\dots\), \(a_m\)
    是 \(1\), \(2\), \(\dots\), \(m-1\) 的排列.
    由假定, 我们可作
    \(
    T_{m-1} = \tau (a_1, \dots, a_{k-1}, a_{k+1}, \dots, a_m)
    \)
    次相邻对换,
    变 \(a_1\), \(\dots\), \(a_{k-1}\),
    \(a_{k+1}\), \(\dots\), \(a_m\)
    为它的自然排列 \(1\), \(2\), \(\dots\), \(m-1\).
    所以, 我们可作 \(T_{m-1} + (m - k)\)~次相邻对换%
    变 \(a_1\), \(a_2\), \(\dots\), \(a_m\)
    为它的自然排列 \(1\), \(2\), \(\dots\), \(m\).

    注意, 每个适合条件 \(1 \leq i < j \leq m\)
    的有序整数对 \((i, j)\)
    恰适合以下三个条件的一个:
    (a)
    \(i \neq k\) 且 \(j \neq k\);
    (b)
    \(i = k\) 且 \(k < j \leq m\);
    (c)
    \(1 \leq i < k\) 且 \(j = k\).
    再注意, 对每个不等于 \(k\),
    且不超过 \(m\) 的正整数 \(\ell\),
    必有 \(a_\ell < m\).
    故
    \begin{align*}
             & \tau (a_1, a_2, \dots, a_m)
        \\
        = {} &
        \sum_{1 \leq i < j \leq m} {\rho(a_i, a_j)}
        \\
        = {} &
        \sum_{\substack{1 \leq i < j \leq m \\i,j \neq k}}
        {\rho(a_i, a_j)}
        + \sum_{k < j \leq m} {\rho(a_k, a_j)}
        + \sum_{1 \leq i < k} {\rho(a_i, a_k)}
        \\
        = {} &
        \tau (a_1, \dots, a_{k-1}, a_{k+1}, \dots, a_m)
        + \sum_{k < j \leq m} {1}
        + \sum_{1 \leq i < k} {0}
        \\
        = {} &
        T_{m-1} + (m - k).
    \end{align*}
    从而, 我们可作
    \(\tau (a_1, a_2, \dots, a_m)\)~次相邻对换%
    变 \(a_1\), \(a_2\), \(\dots\), \(a_m\)
    为它的自然排列 \(1\), \(2\), \(\dots\), \(m\).

    所以, \(P(m)\) 是对的.
    由数学归纳法,
    对任何正整数 \(n\), \(P(n)\) 成立.

    最后, 我们看一般的情形.

    设 \(a_1\), \(a_2\), \(\dots\), \(a_n\)
    的自然排列是 \(c_1\), \(c_2\), \(\dots\), \(c_n\).
    我们记 \(f(c_i) = i\)
    (\(i = 1\), \(2\), \(\dots\), \(n\)).
    (所以, 在 \(a_1\), \(a_2\), \(\dots\), \(a_n\) 中,
    \(a_i\) 是第~\(f(a_i)\) 小的数.)
    那么, 不难看出,
    若 \(d_1\), \(d_2\), \(\dots\), \(d_n\)
    是 \(c_1\), \(c_2\), \(\dots\), \(c_n\)
    的一个排列,
    则 \(f(d_1)\), \(f(d_2)\), \(\dots\), \(f(d_n)\)
    是 \(1\), \(2\), \(\dots\), \(n\)
    的一个排列;
    反过来, 若 \(v_1\), \(v_2\), \(\dots\), \(v_n\) 是
    \(1\), \(2\), \(\dots\), \(n\) 的一个排列,
    则 \(c_{v_1}\), \(c_{v_2}\), \(\dots\), \(c_{v_n}\)
    是 \(c_1\), \(c_2\), \(\dots\), \(c_n\) 的一个排列.
    也不难看出, 若 \(d_t < d_v\), 必有 \(f(d_t) < f(d_v)\).
    反过来, 若 \(f(d_t) < f(d_v)\), 必有 \(d_t < d_v\).
    从而, \(f(a_1)\), \(f(a_2)\), \(\dots\), \(f(a_n)\)
    的逆序数等于 \(a_1\), \(a_2\), \(\dots\), \(a_n\)
    的逆序数.
    我们分别为 \(a_1\), \(a_2\), \(\dots\), \(a_n\)
    取别名 \(f(a_1)\), \(f(a_2)\), \(\dots\), \(f(a_n)\),
    即可用老方法,
    作跟 \(a_1\), \(a_2\), \(\dots\), \(a_n\)
    的逆序数一样多的次数的相邻对换,
    变 \(a_1\), \(a_2\), \(\dots\), \(a_n\) 为它的%
    自然排列 \(c_1\), \(c_2\), \(\dots\), \(c_n\).
    (注意, 我们总可用类似的方法,
    变一般的排列的研究为
    \(1\), \(2\), \(\dots\), \(n\) 的排列的研究.
    所以, 当我们得到
    \(1\), \(2\), \(\dots\), \(n\)
    的排列的只跟逆序有关的性质后,
    我们总可译它为一般的排列的性质.)
\end{proof}

\begin{theorem}
    设排列~I: \(a_1\), \(a_2\), \(\dots\), \(a_n\).
    施一次对换于排列~I,
    得排列~II: \(b_1\), \(b_2\), \(\dots\), \(b_n\).
    则
    \begin{align*}
        s(b_1, b_2, \dots, b_n) = -s(a_1, a_2, \dots, a_n).
    \end{align*}
    (通俗地, 一次对换改变排列的符号.)
\end{theorem}

\begin{proof}
    无妨设排列~I, II 都是
    \(1\), \(2\), \(\dots\), \(n\) 的排列;
    可用我在前面提到的方法推广此事于任何排列.

    设 \(n\)~级单位阵
    \(I = [e_1, e_2, \dots, e_n]\).
    作二个 \(n\)~级阵
    \(A = [e_{a_1}, e_{a_2}, \dots, e_{a_n}]\),
    \(B = [e_{b_1}, e_{b_2}, \dots, e_{b_n}]\).
    那么, \(B\) 可被认为是交换了 \(A\)~的二列,
    且不改变其他的列得到的阵.
    从而,
    \begin{align*}
        %      & s(b_1, b_2, \dots, b_n)
        % \\
        s(b_1, b_2, \dots, b_n)
        = {} & s(b_1, b_2, \dots, b_n) \det {(I)}
        \\
        = {} & \det {(B)}
        \\
        = {} & {-} \det {(A)}
        \\
        = {} & {-} s(a_1, a_2, \dots, a_n) \det {(I)}
        \\
        = {} & {-} s(a_1, a_2, \dots, a_n).
        \qedhere
    \end{align*}
\end{proof}

\begin{theorem}
    设 \(n \geq 2\).
    设 \(a_1\), \(a_2\), \(\dots\), \(a_n\) 是
    \(n\)~个互不相同的整数.
    那么, 在
    \(a_1\), \(a_2\), \(\dots\), \(a_n\) 的全部排列中,
    符号为 \(1\) 的排列%
    跟符号为 \(-1\) 的排列%
    的数目相等.
\end{theorem}

\begin{proof}
    设 \(a_1\), \(a_2\), \(\dots\), \(a_n\)
    的自然排列是 \(c_1\), \(c_2\), \(\dots\), \(c_n\).
    不难看出,
    每个 \(a_1\), \(a_2\), \(\dots\), \(a_n\) 的排列%
    都是 \(c_1\), \(c_2\), \(\dots\), \(c_n\) 的排列,
    且每个 \(c_1\), \(c_2\), \(\dots\), \(c_n\) 的排列%
    都是 \(a_1\), \(a_2\), \(\dots\), \(a_n\) 的排列,

    无妨设 \(c_1\), \(c_2\), \(\dots\), \(c_n\)
    分别为 \(1\), \(2\), \(\dots\), \(n\);
    可用我在前面提到的方法推广此事于任何排列.

    作 \(n\)~级阵 \(U\),
    其中, \([U]_{i,j} = 1\)
    (通俗地, \(U\) 的每个元都是 \(1\)).
    根据交错性, \(\det {(U)} = 0\).
    根据完全展开,
    % (取 \(j_1\), \(j_2\), \(\dots\), \(j_n\)
    % 为 \(1\), \(2\), \(\dots\), \(n\)),
    \begin{align*}
        0
        = {} & \det {(U)}
        \\
        = {} & \sum_{\substack{
        1 \leq i_1, i_2, \dots, i_n \leq n \\
                i_1, i_2, \dots, i_n\,\text{互不相同}
            }}
        {s(i_1, i_2, \dots, i_n)\,
            [U]_{i_1,1} [U]_{i_2,2} \dots [U]_{i_n,n}}
        \\
        = {} & \sum_{\substack{
        1 \leq i_1, i_2, \dots, i_n \leq n \\
                i_1, i_2, \dots, i_n\,\text{互不相同}
            }}
        {s(i_1, i_2, \dots, i_n)}.
    \end{align*}
    我问一个简单的问题:
    设有 \(k\)~个数, 且每个都是 \(1\) 或 \(-1\).
    若它们的和为 \(0\),
    有几个 \(1\), 且有几个 \(-1\)?
    我想, 您能答出,
    \(1\) 跟 \(-1\) 的个数都是 \(k/2\).
\end{proof}

由此可见,
完全展开 \(n\)~级阵 (\(n \geq 2\)) 的行列式的公式里,
符号为 \(+\) 的项的数目%
与符号为 \(-\) 的项的数目%
都是
\((n \cdot (n - 1) \cdot \dots \cdot 2 \cdot 1)/2\).

\vspace{2ex}

% 最后, 我想说公理定义的事.

或许, 您还记得,
我在第一章, 节~\sekcio{14},
``用行列式的性质确定行列式'' 里说过,
行列式的规范性、多线性、交错性可确定行列式:

\begin{theorem}
    设定义在全体 \(n\)~级阵上的函数 \(f\) 适合:

    (1)
    (多线性)
    对任何不超过 \(n\) 的正整数 \(j\),
    任何 \(n-1\)~个 \(n \times 1\)~阵
    \(a_1\), \(\dots\), \(a_{j-1}\),
    \(a_{j+1}\), \(\dots\), \(a_n\),
    任何二个 \(n \times 1\)~阵 \(x\), \(y\),
    任何二个数 \(s\), \(t\),
    有
    \begin{align*}
             & f
            {([a_1, \dots, a_{j-1}, sx + ty,
                        a_{j+1}, \dots, a_n])}
        \\
        = {} &
        s
        f {([a_1, \dots, a_{j-1}, x, a_{j+1}, \dots, a_n])}
        +
        t
        f {([a_1, \dots, a_{j-1}, y, a_{j+1}, \dots, a_n])}.
    \end{align*}

    (2)
    (交错性)
    若 \(n\)~级阵 \(A\) 有二列相同,
    则 \(f {(A)} = 0\).

    那么, 对任何 \(n\)~级阵 \(A\),
    \(f(A) = f(I) \det {(A)}\).

    特别地, 若 \(f(I) = 1\) (规范性),
    则 \(f\) 就是行列式.
\end{theorem}

还记得我如何证此事吗?
我作了一个辅助函数
\begin{align*}
    g(A) = f(A) - f(I) \det {(A)}.
\end{align*}
可以验证, \(g\) 是多线性的、交错性的,
且 \(g(I) = 0\).
于是, 设法证对任何 \(n\)~级阵 \(A\),
\(g(A) = 0\) 即可.
我为什么要作这个 \(g\)?

为解释此事, 我们看,
若我不作此 \(g\), 会如何.
首先, 因为 \(f\) 的多线性与交错性,
我们可类似地写出
\begin{align*}
    f(A)
    = \sum_{\substack{
    1 \leq i_1, i_2, \dots, i_n \leq n \\
            i_1, i_2, \dots, i_n\,\text{互不相同}
        }}
    {[A]_{i_1,1} [A]_{i_2,2} \dots [A]_{i_n,n}\,
        f([e_{i_1}, e_{i_2}, \dots, e_{i_n}])},
\end{align*}
其中, \(e_k\) 是 \(n\)~级单位阵~\(I\) 的列~\(k\).

然后, 我们要确定
\(f([e_{i_1}, e_{i_2}, \dots, e_{i_n}])\).
由多线性与交错性, 我们可推出反称性.
则我们可以适当地交换
\([e_{i_1}, e_{i_2}, \dots, e_{i_n}]\)~的列,
变它为 \([e_1, e_2, \dots, e_n]\),
即 \(n\)~级单位阵 \(I\).
从而, 存在一个由 \(i_1\), \(i_2\), \(\dots\), \(i_n\)
确定的数 \(\sigma (i_1, i_2, \dots, i_n)\),
其等于 \(1\) 或 \(-1\),
使
\begin{align*}
    f([e_{i_1}, e_{i_2}, \dots, e_{i_n}])
    = \sigma (i_1, i_2, \dots, i_n)\, f(I).
\end{align*}
故
\begin{align*}
    f(A)
    = f(I) \sum_{\substack{
    1 \leq i_1, i_2, \dots, i_n \leq n \\
            i_1, i_2, \dots, i_n\,\text{互不相同}
        }}
    {\sigma (i_1, i_2, \dots, i_n)\,
        [A]_{i_1,1} [A]_{i_2,2} \dots [A]_{i_n,n}}.
\end{align*}
问题来了:
\(\sigma (i_1, i_2, \dots, i_n)\) 是什么?

我们回想,
任给 \(1\), \(2\), \(\dots\), \(n\)
的一个排列 \(i_1\), \(i_2\), \(\dots\), \(i_n\),
我们总可作 \(\tau (i_1, i_2, \dots, i_n)\)~次%
相邻对换变它为自然排列.
所以, 我们可作
\(\tau (i_1, i_2, \dots, i_n)\)~次列的交换
(每次交换二列),
变 \([e_{i_1}, e_{i_2}, \dots, e_{i_n}]\) 为 \(I\).
由反称性,
\begin{align*}
    %  &
    f([e_{i_1}, e_{i_2}, \dots, e_{i_n}])
    % \\
    = {} &
    (-1)^{\tau (i_1, i_2, \dots, i_n)}
    f([e_1, e_2, \dots, e_n])
    \\
    = {} &
    s(i_1, i_2, \dots, i_n)\, f(I).
\end{align*}
所以,
% 神秘的
\(\sigma (i_1, i_2, \dots, i_n)
= s(i_1, i_2, \dots, i_n)\).
% 实为老朋友
% 其实就是
% \(s(i_1, i_2, \dots, i_n)\).
再根据完全展开
\begin{align*}
    \det {(A)}
    = {} & \sum_{\substack{
    1 \leq i_1, i_2, \dots, i_n \leq n \\
            i_1, i_2, \dots, i_n\,\text{互不相同}
        }}
    {s(i_1, i_2, \dots, i_n)\,
        [A]_{i_1,1} [A]_{i_2,2} \dots [A]_{i_n,n}},
\end{align*}
% (取 \(j_1\), \(j_2\), \(\dots\), \(j_n\)
% 为 \(1\), \(2\), \(\dots\), \(n\)),
知 \(f(A) = f(I) \det {(A)}\).

由此可见, 理论地, 不作辅助函数 \(g\) 也是可证此事的.
不过, 跟作辅助函数 \(g\) 的论证相比,
这个试直接计算 \(f(A)\) 的论证至少用到了二件事:

(1)
作跟逆序数一样多的次数的相邻对换,
可变一个排列为它的自然排列;

(2)
完全展开行列式的公式.

注意, (1) 是我在本节提到的,
而 (2) (在我的教学里) 是选学的内容.
所以, 为了使论证更好地被理解,
也为了使论证简单,
我作了辅助函数.
\(g\) 的一个特点是 \(g(I) = 0\).
所以, 利用反称性%
变 \(g([e_{i_1}, e_{i_2}, \dots, e_{i_n}])\)
为 \(g(I)\) 时,
我们既不必关注变了几次号,
也不必关注如何变
\(i_1\), \(i_2\), \(\dots\), \(i_n\)
为 \(1\), \(2\), \(\dots\), \(n\)
(比如, 要变 \(3\), \(2\), \(5\), \(1\), \(4\)
为 \(1\), \(2\), \(3\), \(4\), \(5\),
可以从后向前看,
交换 \(5\), \(4\) 的位置,
再交换 \(4\), \(1\) 的位置,
再交换 \(3\), \(1\) 的位置);
我们知道, 可适当地交换文字的位置,
变 \(i_1\), \(i_2\), \(\dots\), \(i_n\)
为 \(1\), \(2\), \(\dots\), \(n\),
即可.
``\(0\) 跟任何数的积都是 \(0\)''
起了大的作用.

用同样的方法, 可证

\begin{theorem}
    设 \(c\) 是常数.
    设定义在全体 \(n\)~级阵上的函数 \(f\) 适合:

    (1)
    (多线性)
    对任何不超过 \(n\) 的正整数 \(j\),
    任何 \(n-1\)~个 \(n \times 1\)~阵
    \(a_1\), \(\dots\), \(a_{j-1}\),
    \(a_{j+1}\), \(\dots\), \(a_n\),
    任何二个 \(n \times 1\)~阵 \(x\), \(y\),
    任何二个数 \(s\), \(t\),
    有
    \begin{align*}
             & f
            {([a_1, \dots, a_{j-1}, sx + ty,
                        a_{j+1}, \dots, a_n])}
        \\
        = {} &
        s
        f {([a_1, \dots, a_{j-1}, x, a_{j+1}, \dots, a_n])}
        +
        t
        f {([a_1, \dots, a_{j-1}, y, a_{j+1}, \dots, a_n])}.
    \end{align*}

    (2)
    (交错性)
    若 \(n\)~级阵 \(A\) 有二列相同,
    则 \(f {(A)} = 0\).

    (3)
    % (``规范性'')
    \(f(I) = c\).

    设定义在全体 \(n\)~级阵上的函数 \(h\)
    也适合这三条性质.
    则对任何 \(n\)~级阵 \(A\), \(f(A) = h(A)\).
\end{theorem}

\begin{proof}
    作辅助函数 \(g(A) = f(A) - h(A)\).
    此 \(g\) 也有多线性与交错性,
    且 \(g(I) = 0\).
    然后, 用我 (在第一章, 节~\sekcio{14}) 用过的方法即可.
    (注意, 此事未提到 ``行列式'' 或 ``det''.)
\end{proof}

\section{Binet--Cauchy 公式}

我要介绍证 Binet--Cauchy 公式的另一个方法.

在第一章, 节~\malneprasekcio{30} 里,
我用完全展开行列式的公式证明了它.
这个证法:

(1)
用简单的记号;

(2)
使您方便地对比选学内容
``Binet--Cauchy 公式''
与必学内容
``Binet--Cauchy 公式 (青春版)''.
% 体会论证的异同
% (其实, 它们大同小异).

当然, 此证法的一个大的缺点是%
用到了完全展开.
% 用到了完全展开
% (或, 组合定义相关的公式).
% 我讲课时, 会尽可能地不跳过 ``选学内容'';
% 不过,
若您不会 (或不想) 用完全展开,
且您想了解如何证明 Binet--Cauchy 公式,
我在这儿给一个利用按多列展开行列式的公式的证明.
至少, 按多列展开行列式 (与它的论证) 不依赖%
排列、逆序数、符号.

这个论证是我初学 Binet--Cauchy 公式时学到的证明.
这个论证体现了重要思想, 算二次;
聪明的数学人利用此它, 得到了多的有名的等式.

为方便, 请允许我引用%
按多列展开行列式公式与 Binet--Cauchy 公式.

\begin{theorem}
    设 \(A\) 是 \(n\)~级阵 (\(n \geq 1\)).
    设 \(k\) 是不超过 \(n\) 的正整数.
    设 \(j_1\), \(j_2\), \(\dots\), \(j_k\) 是%
    不超过 \(n\) 的正整数,
    且 \(j_1 < j_2 < \dots < j_k\).
    则
    \begin{align*}
         &
        \det {(A)}
        = \sum_{1 \leq i_1 < i_2 < \dots < i_k \leq n}
        {\det {\left(
                A\binom{i_1, i_2, \dots, i_k}
                {j_1, j_2, \dots, j_k}
                \right)}}
        \\
         &
        \qquad \qquad \qquad
        \cdot (-1)^{i_1 + i_2 + \dots + i_k
            + j_1 + j_2 + \dots + j_k}
        \det {(A({i_1,i_2,\dots,i_k}|{j_1,j_2,\dots,j_k}))}.
    \end{align*}
\end{theorem}

\begin{theorem}[Binet--Cauchy 公式]
    设 \(A\), \(B\) 分别是 \(m \times n\), \(n \times m\)~阵.

    (1)
    若 \(n > m\), 则 \(\det {(BA)} = 0\).

    (2)
    若 \(n \leq m\),
    则
    \begin{align*}
             & \det {(BA)}
        \\
        = {} & \sum_{1 \leq j_1 < j_2 < \dots < j_n \leq m}
        {
            \det {\left(
                B\binom{1, 2, \dots, n}{j_1, j_2, \dots, j_n}
                \right)}
            \det {\left(
                A\binom{j_1, j_2, \dots, j_n}{1, 2, \dots, n}
                \right)}
        }.
    \end{align*}
    特别地, 若 \(n = m\), 则因适合条件
    \(1 \leq j_1 < j_2 < \dots < j_n \leq m\)
    的 \(j_1\), \(j_2\), \(\dots\), \(j_n\),
    是, 且只能是,
    \(1\), \(2\), \(\dots\), \(n\),
    故
    \begin{align*}
        \det {(BA)}
        = {} &
        \det {\left(
            B\binom{1, 2, \dots, n}{1, 2, \dots, n}
            \right)}
        \det {\left(
            A\binom{1, 2, \dots, n}{1, 2, \dots, n}
            \right)}
        \\
        = {} & \det {(B)} \det {(A)}.
    \end{align*}
\end{theorem}

有几件事值得提.

(1)
若一个方阵有一行的元全为 \(0\),
则它的行列式为 \(0\).

可以直接用定义 (按列~\(1\) 展开行列式) 论证此事;
可用按一行展开行列式的公式论证此事;
当然, 也可利用 (关于行的) 多线性论证此事.
我就不在这儿证了.

(2)
若加一个阵的一列的倍于另一列, 则行列式不变.

这就是我在前面提到的行列式的一个性质.

(3)
设 \(A\), \(B\), \(C\), \(D\)
分别是 \(m \times s\), \(n \times s\),
\(m \times t\), \(n \times t\)~阵.
那么,
\begin{align*}
    \begin{bmatrix}
        A & C \\
        B & D \\
    \end{bmatrix}
\end{align*}
表示一个 \((m + n) \times (s + t)\)-阵 \(M\),
且对任何不超过 \(m + n\)~的正整数 \(i\)
与不超过 \(s + t\)~的正整数 \(j\),
\begin{align*}
    [M]_{i,j}
    = \begin{cases}
          [A]_{i,j},
           & \text{\(i \leq m\), 且 \(j \leq s\)}; \\
          [B]_{i-m,j},
           & \text{\(i > m\), 且 \(j \leq s\)};    \\
          [C]_{i,j-s},
           & \text{\(i \leq m\), 且 \(j > s\)};    \\
          [D]_{i-m,j-s},
           & \text{\(i > m\), 且 \(j > s\)}.
      \end{cases}
\end{align*}

这是我们在第一章, 节~\sekcio{31} 里见过的记号.
现在, 我要进一步地发展此记号.
具体地,
设 \(A\), \(B\)
分别是 \(m \times s\), \(n \times s\)~阵.
那么,
\begin{align*}
    \begin{bmatrix}
        A \\
        B \\
    \end{bmatrix}
\end{align*}
表示一个 \((m + n) \times s\)-阵 \(L\),
且对任何不超过 \(m + n\)~的正整数 \(i\)
与不超过 \(s\)~的正整数 \(j\),
\begin{align*}
    [L]_{i,j}
    = \begin{cases}
          [A]_{i,j},
           & \text{\(i \leq m\)}; \\
          [B]_{i-m,j},
           & \text{\(i > m\)}.
      \end{cases}
\end{align*}

类似地, 设 \(A\), \(C\)
分别是 \(m \times s\), \(m \times t\)~阵.
那么,
\begin{align*}
    \begin{bmatrix}
        A & C \\
    \end{bmatrix}
\end{align*}
表示一个 \(m \times (s + t)\)-阵 \(H\),
且对任何不超过 \(m\)~的正整数 \(i\)
与不超过 \(s + t\)~的正整数 \(j\),
\begin{align*}
    [H]_{i,j}
    = \begin{cases}
          [A]_{i,j},
           & \text{\(j \leq s\)}; \\
          [C]_{i,j-s},
           & \text{\(j > s\)}.
      \end{cases}
\end{align*}

\begin{proof}
    设 \(A\), \(B\) 分别是
    \(m \times n\), \(n \times m\)~阵.
    作 \(m + n\)~级阵
    \begin{align*}
        J =
        \begin{bmatrix}
            I_m & A \\
            -B  & 0 \\
        \end{bmatrix},
    \end{align*}
    其中, 左上角的 \(I_m\) 是 \(m\)~级单位阵,
    且右下角的 \(0\) 是 \(n \times n\)~零阵
    (元全为 \(0\) 的阵).

    我们用二个方式计算 \(J\) 的行列式.

    一方面, 我们按%
    列~\(m + 1\), \(m + 2\), \(\dots\), \(m + n\) 展开,
    有
    \begin{align*}
         &
        \det {(J)}
        = \sum_{1 \leq i_1 < i_2 < \dots < i_n \leq m+n}
        {\det {\left(
                J\binom{i_1,i_2,\dots,i_n}{m+1,m+2,\dots,m+n}
                \right)}}
        \\
         &
        \qquad \qquad \qquad
        \cdot (-1)^{i_1 + i_2 + \dots + i_n
            + (m+1) + (m+2) + \dots + (m+n)}
        \\
         &
        \qquad \qquad \qquad
        \cdot \det {(J({i_1,i_2,\dots,i_n}|{m+1,m+2,\dots,m+n}))}.
    \end{align*}
    注意, 若 \(i_n > m\), 则因
    \(\displaystyle
    J\binom{i_1,i_2,\dots,i_n}{m+1,m+2,\dots,m+n}\)
    的行~\(n\) 的元全为 \(0\),
    故这个子阵的行列式为 \(0\).
    特别地, 若 \(n > m\), 则 \(i_n\) 必高于 \(m\).
    故 \(n > m\) 时, \(J\) 的行列式为 \(0\).

    当 \(n \leq m\) 时,
    我们去除使 \(i_n > m\) 的
    \(i_1\), \(i_2\), \(\dots\), \(i_n\),
    有
    \begin{align*}
         &
        \det {(J)}
        = \sum_{1 \leq i_1 < i_2 < \dots < i_n \leq m}
        {\det {\left(
                J\binom{i_1,i_2,\dots,i_n}{m+1,m+2,\dots,m+n}
                \right)}}
        \\
         &
        \qquad \qquad \qquad
        \cdot (-1)^{i_1 + i_2 + \dots + i_n
            + (m+1) + (m+2) + \dots + (m+n)}
        \\
         &
        \qquad \qquad \qquad
        \cdot \det {(J({i_1,i_2,\dots,i_n}|{m+1,m+2,\dots,m+n}))}.
    \end{align*}
    不难看出,
    因为 \(1 \leq i_1 < i_2 < \dots < i_n \leq m\),
    \begin{align*}
        J\binom{i_1,i_2,\dots,i_n}{m+1,m+2,\dots,m+n}
        = A\binom{i_1,i_2,\dots,i_n}{1,2,\dots,n},
    \end{align*}
    且
    \begin{align*}
        J({i_1,i_2,\dots,i_n}|{m+1,m+2,\dots,m+n})
        =
        \begin{bmatrix}
            I_m (i_1,i_2,\dots,i_n|) \\
            -B
        \end{bmatrix},
    \end{align*}
    其中,
    \(I_m (i_1,i_2,\dots,i_n|)\)
    表示去除 \(I_m\)~的%
    行~\(i_1\), \(i_2\), \(\dots\), \(i_n\)
    后 (不改变列) 得到的阵.
    按列~\(i_1\), \(i_2\), \(\dots\), \(i_n\) 展开
    \(T_{i_1,i_2,\dots,i_n}
    = J({i_1,i_2,\dots,i_n}|{m+1,m+2,\dots,m+n})\)~的%
    行列式, 有
    \begin{align*}
             & \det {(T_{i_1,i_2,\dots,i_n})}
        \\
        = {} &
        \sum_{1 \leq k_1 < k_2 < \dots < k_n \leq m}
        {\det {\left(
                T_{i_1,i_2,\dots,i_n}
                \binom{k_1, k_2, \dots, k_n}{i_1, i_2, \dots, i_n}
                \right)}}
        \\
             &
        \qquad
        \cdot (-1)^{k_1 + k_2 + \dots + k_n
            + i_1 + i_2 + \dots + i_n}
        \det {(T_{i_1,i_2,\dots,i_n}
        ({k_1,k_2,\dots,k_n}|{i_1,i_2,\dots,i_n}))}.
    \end{align*}
    注意, 若 \(k_1 \leq m-n\),
    则因 \(\displaystyle T_{i_1,i_2,\dots,i_n}
    \binom{k_1,k_2,\dots,k_n}{i_1,i_2,\dots,i_n}\)
    的行~\(1\) 的元全为 \(0\),
    故这个子阵的行列式为 \(0\).
    并且, 若 \(k_1 > m-n\),
    则 \(k_1\), \(k_2\), \(\dots\), \(k_n\)
    是, 且只能是
    \(m-n+1\), \(m-n+2\), \(\dots\), \(m\).
    故
    \begin{align*}
             & \det {(T_{i_1,i_2,\dots,i_n})}
        \\
        = {} &
        \hphantom{\cdot\,\,}
        \det {\left(
            T_{i_1,i_2,\dots,i_n}
            \binom{m-n+1,m-n+2,\dots,m}{i_1,i_2,\dots,i_n}
            \right)}
        \\
             &
        \cdot (-1)^{(m-n+1) + (m-n+2) + \dots + m
            + i_1 + i_2 + \dots + i_n}
        \\
             &
        \cdot
        \det {(T_{i_1,i_2,\dots,i_n}
        ({m-n+1,m-n+2,\dots,m}|{i_1,i_2,\dots,i_n}))}
        \\
        = {} &
        \hphantom{\cdot\,\,}
        \det {\left(
            (-B)
            \binom{1,2,\dots,n}{i_1,i_2,\dots,i_n}
            \right)}
        \\
             & \cdot
        (-1)^{(m-n+1) + (m-n+2) + \dots + (m-n+n)
            + i_1 + i_2 + \dots + i_n}
        \\
             & \cdot
        \det {(I_m
        ({i_1,i_2,\dots,i_n}|{i_1,i_2,\dots,i_n}))}
        \\
        = {} &
        (-1)^{(m-n+1) + (m-n+2) + \dots + (m-n+n)
                + i_1 + i_2 + \dots + i_n}\,
        (-1)^n
        \det {\left(
            B
            \binom{1,2,\dots,n}{i_1,i_2,\dots,i_n}
            \right)}.
    \end{align*}
    从而
    \begin{align*}
        \det {(J)}
        = {} & \sum_{1 \leq i_1 < i_2 < \dots < i_n \leq m}
        {\det {\left(
                A\binom{i_1,i_2,\dots,i_n}{1,2,\dots,n}
                \right)}}
        \\
             & \qquad
        \cdot
        (-1)^{i_1 + i_2 + \dots + i_n
            + (m+1) + (m+2) + \dots + (m+n)}
        \\
             & \qquad
        \cdot
        (-1)^{(m-n+1) + (m-n+2) + \dots + (m-n+n)
            + i_1 + i_2 + \dots + i_n}\, (-1)^n
        \\
             & \qquad
        \cdot
        \det {\left(
            B
            \binom{1,2,\dots,n}{i_1,i_2,\dots,i_n}
            \right)}.
    \end{align*}
    注意,
    \begin{align*}
             &
        \hphantom{\cdot\,\,}
        (-1)^{i_1 + i_2 + \dots + i_n
            + (m+1) + (m+2) + \dots + (m+n)}
        \\
             &
        \cdot
        (-1)^{(m-n+1) + (m-n+2) + \dots + (m-n+n)
            + i_1 + i_2 + \dots + i_n}
        \cdot
        (-1)^n
        \\
        = {} & (-1)^{2(i_1+i_2+\dots+i_n)}
        \cdot (-1)^{2mn}
        \cdot (-1)^{(n-(n-1))+((n-1)-(n-2))+\dots+(1-(n-n))}
        \cdot (-1)^{n}
        \\
        = {} & (-1)^n \cdot (-1)^n
        \\
        = {} & 1,
    \end{align*}
    故
    \begin{align*}
        \det {(J)}
        = \sum_{1 \leq i_1 < i_2 < \dots < i_n \leq m}
        {\det {\left(
                B
                \binom{1,2,\dots,n}{i_1,i_2,\dots,i_n}
                \right)}
            \det {\left(
                A\binom{i_1,i_2,\dots,i_n}{1,2,\dots,n}
                \right)}}.
    \end{align*}

    另一方面, 我们也可用行列式的性质,
    施适当的变换于 \(J\),
    得一个新阵 \(K\),
    且 \(J\), \(K\) 的行列式相等.
    具体地, 加
    \(J\)~的列~\(1\) 的 \(-[A]_{1,1}\)~倍于列~\(m+1\),
    \(J\)~的列~\(2\) 的 \(-[A]_{2,1}\)~倍于列~\(m+1\),
    ……
    \(J\)~的列~\(m\) 的 \(-[A]_{m,1}\)~倍于列~\(m+1\),
    得
    \begin{align*}
        J_1 =
        \begin{bmatrix}
            1          & 0          & \cdots & 0          &
            0          & [A]_{1,2}  & \cdots & [A]_{1,n}    \\
            0          & 1          & \cdots & 0          &
            0          & [A]_{2,2}  & \cdots & [A]_{2,n}    \\
            \vdots     & \vdots     & {}     & \vdots     &
            \vdots     & \vdots     & {}     & \vdots       \\
            0          & 0          & \cdots & 1          &
            0          & [A]_{m,2}  & \cdots & [A]_{m,n}    \\
            -[B]_{1,1} & -[B]_{1,2} & \cdots & -[B]_{1,m} &
            p_{1,1}    & 0          & \cdots & 0            \\
            -[B]_{2,1} & -[B]_{2,2} & \cdots & -[B]_{2,m} &
            p_{2,1}    & 0          & \cdots & 0            \\
            \vdots     & \vdots     & {}     & \vdots     &
            \vdots     & \vdots     & {}     & \vdots       \\
            -[B]_{n,1} & -[B]_{n,2} & \cdots & -[B]_{n,m} &
            p_{n,1}    & 0          & \cdots & 0            \\
        \end{bmatrix},
    \end{align*}
    其中,
    \begin{align*}
        p_{k,1}
        = {} &
        (-[B]_{k,1})(-[A]_{1,1})
        + (-[B]_{k,2})(-[A]_{2,1})
        + \dots
        + (-[B]_{k,m})(-[A]_{m,1})
        \\
        = {} &
        [B]_{k,1} [A]_{1,1}
        + [B]_{k,2} [A]_{2,1}
        + \dots
        + [B]_{k,m} [A]_{m,1}
        \\
        = {} &
        [BA]_{k,1}.
    \end{align*}
    根据行列式的性质,
    \(\det {(J)} = \det {(J_1)}\).

    加
    \(J_1\)~的列~\(1\) 的 \(-[A]_{1,2}\)~倍于列~\(m+2\),
    \(J_1\)~的列~\(2\) 的 \(-[A]_{2,2}\)~倍于列~\(m+2\),
    ……
    \(J_1\)~的列~\(m\) 的 \(-[A]_{m,2}\)~倍于列~\(m+2\),
    得
    \begin{align*}
        J_2 =
        \begin{bmatrix}
            1          & 0          & \cdots & 0          &
            0          & 0          & \cdots & [A]_{1,n}    \\
            0          & 1          & \cdots & 0          &
            0          & 0          & \cdots & [A]_{2,n}    \\
            \vdots     & \vdots     & {}     & \vdots     &
            \vdots     & \vdots     & {}     & \vdots       \\
            0          & 0          & \cdots & 1          &
            0          & 0          & \cdots & [A]_{m,n}    \\
            -[B]_{1,1} & -[B]_{1,2} & \cdots & -[B]_{1,m} &
            p_{1,1}    & p_{1,2}    & \cdots & 0            \\
            -[B]_{2,1} & -[B]_{2,2} & \cdots & -[B]_{2,m} &
            p_{2,1}    & p_{2,2}    & \cdots & 0            \\
            \vdots     & \vdots     & {}     & \vdots     &
            \vdots     & \vdots     & {}     & \vdots       \\
            -[B]_{n,1} & -[B]_{n,2} & \cdots & -[B]_{n,m} &
            p_{n,1}    & p_{n,2}    & \cdots & 0            \\
        \end{bmatrix},
    \end{align*}
    其中,
    \begin{align*}
        p_{k,2}
        = {} &
        (-[B]_{k,1})(-[A]_{1,2})
        + (-[B]_{k,2})(-[A]_{2,2})
        + \dots
        + (-[B]_{k,m})(-[A]_{m,2})
        \\
        = {} &
        [B]_{k,1} [A]_{1,2}
        + [B]_{k,2} [A]_{2,2}
        + \dots
        + [B]_{k,m} [A]_{m,2}
        \\
        = {} &
        [BA]_{k,2}.
    \end{align*}
    根据行列式的性质,
    \(\det {(J_1)} = \det {(J_2)}\),
    故
    \(\det {(J)} = \det {(J_2)}\).

    ……

    加
    \(J_{n-1}\)~的列~\(1\) 的 \(-[A]_{1,n}\)~倍于列~\(m+n\),
    \(J_{n-1}\)~的列~\(2\) 的 \(-[A]_{2,n}\)~倍于列~\(m+n\),
    ……
    \(J_{n-1}\)~的列~\(m\) 的 \(-[A]_{m,n}\)~倍于列~\(m+n\),
    得
    \begin{align*}
        J_n
        = {} &
        \begin{bmatrix}
            1          & 0          & \cdots & 0          &
            0          & 0          & \cdots & 0            \\
            0          & 1          & \cdots & 0          &
            0          & 0          & \cdots & 0            \\
            \vdots     & \vdots     & {}     & \vdots     &
            \vdots     & \vdots     & {}     & \vdots       \\
            0          & 0          & \cdots & 1          &
            0          & 0          & \cdots & 0            \\
            -[B]_{1,1} & -[B]_{1,2} & \cdots & -[B]_{1,m} &
            p_{1,1}    & p_{1,2}    & \cdots & p_{1,n}      \\
            -[B]_{2,1} & -[B]_{2,2} & \cdots & -[B]_{2,m} &
            p_{2,1}    & p_{2,2}    & \cdots & p_{2,n}      \\
            \vdots     & \vdots     & {}     & \vdots     &
            \vdots     & \vdots     & {}     & \vdots       \\
            -[B]_{n,1} & -[B]_{n,2} & \cdots & -[B]_{n,m} &
            p_{n,1}    & p_{n,2}    & \cdots & p_{n,n}      \\
        \end{bmatrix}
        \\
        = {} &
        \begin{bmatrix}
            I_m & 0  \\
            -B  & BA \\
        \end{bmatrix},
    \end{align*}
    其中, 右上角的 \(0\) 是 \(m \times n\)~零阵,
    且
    \begin{align*}
        p_{k,n}
        = {} &
        (-[B]_{k,1})(-[A]_{1,n})
        + (-[B]_{k,2})(-[A]_{2,n})
        + \dots
        + (-[B]_{k,m})(-[A]_{m,n})
        \\
        = {} &
        [B]_{k,1} [A]_{1,n}
        + [B]_{k,2} [A]_{2,n}
        + \dots
        + [B]_{k,m} [A]_{m,n}
        \\
        = {} &
        [BA]_{k,n}.
    \end{align*}
    根据行列式的性质,
    \(\det {(J_n)} = \det {(J_{n-1})}\),
    故
    \(\det {(J)} = \det {(J_n)}\).

    我们要如何计算 \(J_n\)~的行列式?
    我给您三个方法:
    (a)
    用第一章, 节~\sekcio{31} 的第~2~个例
    (计算 \(J_n\) 的转置的行列式);
    (b)
    按行~\(1\), \(2\), \(\dots\), \(m\) 展开;
    (c)
    按列~\(m+1\), \(m+2\), \(\dots\), \(m+n\) 展开.
    若您没错, 您可算出,
    \begin{align*}
        \det {(J)}
        = \det {(J_n)}
        = \det {(I_m)} \det {(BA)}
        = \det {(BA)}.
    \end{align*}

    现在, 我们比较二次计算的结果.
    这就是 Binet--Cauchy 公式.
\end{proof}

我还想说一些话.

设 \(A\), \(B\) 分别是 \(s \times n\) 与 \(m \times s\)~阵.
不难看出, \(BA\) 有意义, 且 \(BA\) 是一个 \(m \times n\)~阵.
\(BA\) 可能不是一个方阵, 故说 \(BA\) 的行列式可能是无意义的.
不过, 我们仍可说 \(BA\)~的子阵的行列式, 若这个子阵是方阵.
设 \(k\) 是一个既不超过 \(m\), 也不超过 \(n\) 的正整数.
设 \(1 \leq i_1 < \dots < i_k \leq m\);
设 \(1 \leq \ell_1 < \dots < \ell_k \leq n\).
我们说, 我们可以用 \(B\) 与 \(A\) 的 \(k\)~级子阵的行列式%
写出 \(BA\)~的 \(k\)~级子阵
\begin{align*}
    F = (BA)\binom{i_1,\dots,i_k}{\ell_1,\dots,\ell_k}
\end{align*}
的行列式.

首先, 注意,
\begin{align*}
    (BA)\binom{i_1,\dots,i_k}{\ell_1,\dots,\ell_k}
    =
    B\binom{i_1,\dots,i_k}{1,\dots,s}
    \,
    A\binom{1,\dots,s}{\ell_1,\dots,\ell_k}.
\end{align*}
为说明此事, 我们设
\begin{align*}
    D = B\binom{i_1,\dots,i_k}{1,\dots,s},
    \quad
    C = A\binom{1,\dots,s}{\ell_1,\dots,\ell_k}.
\end{align*}
则 \(D\), \(C\) 分别是 \(k \times s\), \(s \times k\)~阵,
\([D]_{u,j} = [B]_{i_u,j}\),
且 \([C]_{i,v} = [A]_{i,\ell_v}\).
则
\begin{align*}
    [DC]_{u,v}
    = {} &
    \sum_{p = 1}^{s} {[D]_{u,p} [C]_{p,v}}
    \\
    = {} &
    \sum_{p = 1}^{s} {[B]_{i_u,p} [A]_{p,\ell_v}}
    \\
    = {} &
    [BA]_{i_u,\ell_v}
    \\
    = {} &
    [F]_{u,v}.
\end{align*}

现在, 我们可以对 \(C\), \(D\) 用 Binet--Cauchy 公式了.
我再说一次, \(C\) 是 \(s \times k\)~的,
且 \(D\) 是 \(k \times s\)~的.
则

(1)
若 \(k > s\), 则 \(\det {(DC)} = 0\).

(2)
若 \(k \leq s\), 则
\begin{align*}
         & \det {(DC)}
    \\
    = {} & \sum_{1 \leq j_1 < \dots < j_k \leq s}
    {
        \det {\left(
            D\binom{1, \dots, k}{j_1, \dots, j_k}
            \right)}
        \det {\left(
            C\binom{j_1, \dots, j_k}{1, \dots, k}
            \right)}
    }
    \\
    = {} & \sum_{1 \leq j_1 < \dots < j_k \leq s}
    {
        \det {\left(
            B\binom{i_1, \dots, i_k}{j_1, \dots, j_k}
            \right)}
        \det {\left(
            A\binom{j_1, \dots, j_k}{\ell_1, \dots, \ell_k}
            \right)}
    }.
\end{align*}

综上, 我们有

\begin{theorem}
    设 \(A\), \(B\) 分别是 \(s \times n\) 与 \(m \times s\)~阵.
    设 \(k\) 是一个既不超过 \(m\), 也不超过 \(n\) 的正整数.
    设 \(1 \leq i_1 < \dots < i_k \leq m\);
    设 \(1 \leq \ell_1 < \dots < \ell_k \leq n\).

    (1)
    若 \(k > s\), 则
    \begin{align*}
        \det {\left(
            (BA)\binom{i_1,\dots,i_k}{\ell_1,\dots,\ell_k}
            \right)} = 0.
    \end{align*}

    (2)
    若 \(k \leq s\), 则
    \begin{align*}
             & \det {\left(
            (BA)\binom{i_1,\dots,i_k}{\ell_1,\dots,\ell_k}
            \right)}
        \\
        = {} & \sum_{1 \leq j_1 < \dots < j_k \leq s}
        {
            \det {\left(
                B\binom{i_1, \dots, i_k}{j_1, \dots, j_k}
                \right)}
            \det {\left(
                A\binom{j_1, \dots, j_k}{\ell_1, \dots, \ell_k}
                \right)}
        }.
    \end{align*}
\end{theorem}

可视此事为 Binet--Cauchy 公式的一个推广:
若 \(m = n\), \(k = n\),
且
\(i_u = \ell_u = u\) (\(u = 1\), \(2\), \(\dots\), \(n\)),
则这就是 Binet--Cauchy 公式.

\vspace{2ex}

若我们取 \(m = n = s\), 则

\begin{theorem}
    设 \(A\), \(B\) 是 \(n\)~级阵.
    设 \(k\) 是一个不超过 \(n\) 的正整数.
    设 \(1 \leq i_1 < \dots < i_k \leq n\);
    设 \(1 \leq \ell_1 < \dots < \ell_k \leq n\).
    则
    \begin{align*}
             & \det {\left(
            (BA)\binom{i_1,\dots,i_k}{\ell_1,\dots,\ell_k}
            \right)}
        \\
        = {} & \sum_{1 \leq j_1 < \dots < j_k \leq n}
        {
            \det {\left(
                B\binom{i_1, \dots, i_k}{j_1, \dots, j_k}
                \right)}
            \det {\left(
                A\binom{j_1, \dots, j_k}{\ell_1, \dots, \ell_k}
                \right)}
        }.
    \end{align*}
\end{theorem}

我们知道, \(A\) 的一个子阵可被认为是%
取 \(A\) 的于一些行与一些列的元%
按原来的次序作成的新阵,
且也可被认为是%
去除 \(A\) 的一些行与一些列后,
剩下的元按原来的次序作成的阵.
所以, 我们也可如此改写前面的结论:

\begin{theorem}
    设 \(A\), \(B\) 是 \(n\)~级阵.
    设 \(k\) 是一个不超过 \(n\) 的正整数.
    设 \(1 \leq i_1 < \dots < i_k \leq n\);
    设 \(1 \leq \ell_1 < \dots < \ell_k \leq n\).
    则
    \begin{align*}
             & \det {(
        (BA)({i_1,\dots,i_k}|{\ell_1,\dots,\ell_k})
        )}
        \\
        = {} & \sum_{1 \leq j_1 < \dots < j_k \leq n}
        {
        \det {(
        B({i_1, \dots, i_k}|{j_1, \dots, j_k})
        )}
        \det {(
        A({j_1, \dots, j_k}|{\ell_1, \dots, \ell_k})
        )}
        }.
    \end{align*}
\end{theorem}

最后, 我们看一个例.

\begin{example}
    设命题 \(P(n)\):
    对任何 \(n\)~级阵 \(A\), \(B\),
    有 \(\det {(BA)} = \det {(B)} \det {(A)}\).
    我们用数学归纳法证明,
    对任何正整数 \(n\), \(P(n)\) 是对的.

    \(P(1)\) 是对的.
    设 \(A = [a]\), 且 \(B = [b]\).
    则 \(BA = [ba]\).
    则 \(\det {(BA)} = ba = \det {(B)} \det {(A)}\).

    现在, 我们假定 \(P(n-1)\) 是对的.
    我们要证 \(P(n)\) 是对的.
    设 \(A\), \(B\) 是 \(n\)~级阵.
    则
    \begin{align*}
        \det {(BA)}
        = {} &
        \sum_{p=1}^{n} {
        (-1)^{p+1} [BA]_{p,1} \det {((BA)(p|1))}
        }
        \\
        = {} &
        \sum_{p=1}^{n} {
        (-1)^{p+1}
        \left(
        \sum_{q=1}^{n} {
            [B]_{p,q} [A]_{q,1}
        }
        \right)
        \det {((BA)(p|1))}
        }
        \\
        = {} &
        \sum_{p=1}^{n} {
        \sum_{q=1}^{n} {
        (-1)^{p+1}
            [B]_{p,q} [A]_{q,1}
        \det {((BA)(p|1))}
        }
        }
        \\
        = {} &
        \sum_{q=1}^{n} {
        \sum_{p=1}^{n} {
        (-1)^{p+1}
            [B]_{p,q} [A]_{q,1}
        \det {((BA)(p|1))}
        }
        }
        \\
        = {} &
        \sum_{q=1}^{n} {
        [A]_{q,1}
        \sum_{p=1}^{n} {
        (-1)^{p+1}
            [B]_{p,q}
        \det {((BA)(p|1))}
        }
        }.
    \end{align*}
    作 \(n\)~级阵 \(C_q\) 如下:
    \begin{align*}
        [C_q]_{i,j} =
        \begin{cases}
            [B]_{i,q},  & j = 1;    \\
            [BA]_{i,j}, & j \neq 1.
        \end{cases}
    \end{align*}
    通俗地, \(C_q\)~的列~\(1\) 是 \(B\)~的列~\(q\),
    且 \(C_q\)~的列~\(j\) 是 \(BA\)~的列~\(j\)
    (\(j \neq 1\)).
    则 \([B]_{p,q} = [C_q]_{p,1}\),
    且 \((BA)(p|1) = C_q (p|1)\).
    则
    \begin{align*}
        \det {(BA)}
        = {} &
        \sum_{q=1}^{n} {
        [A]_{q,1}
        \sum_{p=1}^{n} {
        (-1)^{p+1}
            [B]_{p,q}
        \det {((BA)(p|1))}
        }
        }
        \\
        = {} &
        \sum_{q=1}^{n} {
        [A]_{q,1}
        \sum_{p=1}^{n} {
        (-1)^{p+1}
            [C_q]_{p,1}
        \det {(C_q (p|1))}
        }
        }
        \\
        = {} &
        \sum_{q=1}^{n} {
        [A]_{q,1}
        \det {(C_q)}
        }.
    \end{align*}
    注意, \(j \neq 1\) 时,
    \begin{align*}
        [C_q]_{i,j}
        = {} &
        [BA]_{i,j}
        \\
        = {} &
        \sum_{k=1}^{n} {
        [B]_{i,k} [A]_{k,j}
        }
        \\
        = {} &
        [B]_{i,q} [A]_{q,j} +
        \sum_{\substack{1 \leq k \leq n \\k \neq q}} {
            [B]_{i,k} [A]_{k,j}
        }
        \\
        = {} &
        [B]_{i,q} [A]_{q,j} +
        \sum_{\ell=1}^{n-1} {
        [B]_{i,\ell+\rho(\ell,k)}
            [A]_{\ell+\rho(\ell,k),j}
        }
        \\
        = {} &
        [B]_{i,q} [A]_{q,j} +
        \sum_{\ell=1}^{n-1} {
        [B({|q})]_{i,\ell}
            [A({q|})]_{\ell,j}
        }
        \\
        = {} &
        [B]_{i,q} [A]_{q,j} +
        [B({|q})\, A({q|})]_{i,j},
    \end{align*}
    其中, \(B({|q})\) 是去除 \(B\)~的列~\(q\) 后 (不改变行)
    得到的 \(n \times (n-1)\)~阵,
    且 \(A({q|})\) 是去除 \(A\)~的行~\(q\) 后 (不改变列)
    得到的 \((n-1) \times n\)~阵.
    加
    \(C_q\)~的列~\(1\) 的 \(-[A]_{q,2}\)~倍于列~\(2\),
    \(C_q\)~的列~\(1\) 的 \(-[A]_{q,3}\)~倍于列~\(3\),
    ……
    \(C_q\)~的列~\(1\) 的 \(-[A]_{q,n}\)~倍于列~\(n\),
    得 \(n\)~级阵 \(D_q\).
    则
    \begin{align*}
        [D_q]_{i,j} =
        \begin{cases}
            [B]_{i,q},                 & j = 1;    \\
            [B({|q})\, A({q|})]_{i,j}, & j \neq 1.
        \end{cases}
    \end{align*}
    根据行列式的性质, \(\det {(C_q)} = \det {(D_q)}\).
    注意, \([D_q]_{p,1} = [B]_{p,q}\),
    且 \(D_q (p|1) = (B({|q})\, A({q|}))\, (p|1)\).
    则
    \begin{align*}
        \det {(BA)}
        = {} &
        \sum_{q=1}^{n} {
        [A]_{q,1}
        \det {(C_q)}
        }
        \\
        = {} &
        \sum_{q=1}^{n} {
        [A]_{q,1}
        \det {(D_q)}
        }
        \\
        = {} &
        \sum_{q=1}^{n} {
        [A]_{q,1}
        \sum_{p=1}^{n} {
        (-1)^{p+1}
            [D_q]_{p,1}
        \det {(D_q (p|1))}
        }
        }
        \\
        = {} &
        \sum_{q=1}^{n} {
        [A]_{q,1}
        \sum_{p=1}^{n} {
        (-1)^{p+1}
            [B]_{p,q}
        \det {((B({|q})\, A({q|}))\, (p|1))}
        }
        }.
    \end{align*}
    注意, 对 \(m \times s\)~阵 \(X\)
    与 \(s \times n\)~阵 \(Y\),
    对正整数 \(u \leq m\)
    与正整数 \(v \leq n\),
    有 \((XY)(u|v) = X({u|})\, Y({|v})\).
    则
    \begin{align*}
        (B({|q})\, A({q|}))\, (p|1)
        = (B({|q}))({p|})\, (A({q|}))({|1})
        = B({p|q})\, A({q|1}).
    \end{align*}
    于是, 由假定,
    \begin{align*}
        \det {((B({|q})\, A({q|}))\, (p|1))}
        = \det {(B({p|q})\, A({q|1}))}
        = \det {(B({p|q}))} \det {(A({q|1}))}.
    \end{align*}
    则
    \begin{align*}
        \det {(BA)}
        = {} &
        \sum_{q=1}^{n} {
        [A]_{q,1}
        \sum_{p=1}^{n} {
        (-1)^{p+1}
            [B]_{p,q}
        \det {((B({|q})\, A({q|}))\, (p|1))}
        }
        }
        \\
        = {} &
        \sum_{q=1}^{n} {
        [A]_{q,1}
        \sum_{p=1}^{n} {
        (-1)^{p+1}
            [B]_{p,q}
        \det {(B({p|q})\, A({q|1}))}
        }
        }
        \\
        = {} &
        \sum_{q=1}^{n} {
        [A]_{q,1}
        \sum_{p=1}^{n} {
        (-1)^{p+1}
            [B]_{p,q}
        \det {(B({p|q}))} \det {(A({q|1}))}
        }
        }
        \\
        = {} &
        \sum_{q=1}^{n} {
        [A]_{q,1}
        \sum_{p=1}^{n} {
        (-1)^{p+q} (-1)^{q+1}
            [B]_{p,q}
        \det {(B({p|q}))} \det {(A({q|1}))}
        }
        }
        \\
        = {} &
        \sum_{q=1}^{n} {
        [A]_{q,1}
        (-1)^{q+1}
        \det {(A({q|1}))}
        \sum_{p=1}^{n} {
        (-1)^{p+q}
            [B]_{p,q}
        \det {(B({p|q}))}
        }
        }
        \\
        = {} &
        \sum_{q=1}^{n} {
        [A]_{q,1}
        (-1)^{q+1}
        \det {(A({q|1}))}
        \det {(B)}
        }
        \\
        = {} &
        \det {(B)}
        \sum_{q=1}^{n} {
        [A]_{q,1}
        (-1)^{q+1}
        \det {(A({q|1}))}
        }
        \\
        = {} &
        \det {(B)} \det {(A)}.
    \end{align*}
    所以, \(P(n)\) 是对的.
    由数学归纳法, 待证命题成立.
\end{example}

\end{document}

This material is not very interesting.

Some properties of determinants are given here
(along with some examples of
finding determinants of some famous square matrices);
some properties of permutations are given here
(along with some remarks on
how to determine the determinant function);
another proof of the formula of Binet--Cauchy
is given here
(along with some generalisations of it).
